\section*{الفصل \ref{ch:b3:electrode-kinetics} --- حركية الأقطاب والتحليل الكهربائي}

\begin{solution}{exo:b3:electrode-kinetics:1}
الميل المهبطي $2.303RT/\alpha F = \qty{118}{mV}$ يعطي $\alpha = 59.2/118 = 0.50$؛ و$j_0 = \qty{2.0e-6}{A.cm^{-2}}$،
وهو التقاطع عند $\eta = 0$.
\end{solution}

\begin{solution}{exo:b3:electrode-kinetics:2}
$i_0 = \qty{2.0e-4}{A}$؛ $R_{\mathrm{ct}} = RT/(Fi_0) = 0.02569/\num{2.0e-4} = \qty{128}{\ohm}$.
\end{solution}

\begin{solution}{exo:b3:electrode-kinetics:3}
من أجل \ce{MgSO4}، $I = \frac12(4c + 4c) = 4c$. عند \qty{1.0}{mmol/L}، $I = \qty{4.0}{mol.m^{-3}}$ و$\kappa^{-1} =
\qty{0.96}{nm} \times \sqrt{100/4} = \qty{4.8}{nm}$؛ وعند \qty{0.10}{mol/L}، $I = \qty{400}{mol.m^{-3}}$، $\kappa^{-1} = \qty{0.48}{nm}$.
\end{solution}

\begin{solution}{exo:b3:electrode-kinetics:4}
$i_{\mathrm c} = C_{\mathrm{dl}}Av = \num{20e-6} \times 0.0707 \times 0.1 = \qty{0.14}{\micro A}$، أي أقل من 1\,\% من القمة. وعند
\qty{10}{V/s} يبلغ \qty{14}{\micro A}، بينما لا تكبر القمة إلا إلى $19\sqrt{100} = \qty{190}{\micro A}$:
فتيار الشحن يكبر وفق $v$، والتيار الفارادي وفق $\sqrt v$.
\end{solution}

\begin{solution}{exo:b3:electrode-kinetics:5}
$D = \pi\bigl(i\sqrt t/nFAc^*\bigr)^2 = \pi(\num{12.2e-6}/(96\,485 \times 0.0707 \times \num{1.00e-6}))^2 =
\qty{1.0e-5}{cm^2.s^{-1}}$.
\end{solution}

\begin{solution}{exo:b3:electrode-kinetics:6}
$FAc^* = \qty{1.364e-2}{C.cm^{-1}}$، $F/RT = \qty{38.92}{V^{-1}}$. عند \qty{50}{mV/s}: $\sqrt{38.92 \times 0.05 \times \num{7.0e-6}} =
\num{3.69e-3}$، $i_{\mathrm p} = 0.4463 \times \num{1.364e-2} \times \num{3.69e-3} = \qty{22.5}{\micro A}$؛ وعند \qty{500}{mV/s}، أكبر بمقدار $\sqrt{10}$
مرة، أي \qty{71}{\micro A}.
\end{solution}

\begin{solution}{exo:b3:electrode-kinetics:7}
(أ) عكوسة. (ب) \omterm{def:b3:electrode-kinetics:reversibility}{شبه عكوسة}: المسافة بين القمتين تكبر مع سرعة المسح.
(ج) تعقب انتقالَ الإلكترون خطوةٌ كيميائية (EC): ففي المسوح البطيئة يُستهلك
الناتج قبل المسح العكسي؛ أما المسح السريع فيسبق الكيمياء.
\end{solution}

\begin{solution}{exo:b3:electrode-kinetics:8}
$K_{ij}a_j/a_i = \num{2e-4} \times 0.14/\num{4.0e-3} = 0.007$: فالقطب يقرأ قيمة أعلى بنسبة 0.7\,\%.
\end{solution}

\begin{solution}{exo:b3:electrode-kinetics:9}
المربعات الصغرى: $b = \qty{0.1093}{\micro A.L.\micro g^{-1}}$، $a = \qty{0.466}{\micro A}$، $s = \qty{0.011}{\micro A}$؛ $c_0 = a/b =
\qty{4.26}{\micro g.L^{-1}}$، مع $u(c_0) = \frac sb\sqrt{\frac1n + \frac{\bar y^2}{b^2S_{xx}}} = \qty{0.12}{\micro g.L^{-1}}$
($n = 4$، $\bar y = 1.286$، $S_{xx} = 125$).
\end{solution}

\begin{solution}{exo:b3:electrode-kinetics:10}
$Q = 0.0500 \times 1800 = \qty{90.0}{C}$؛ $n(\ce{Cu}) = Q/2F = \qty{4.66e-4}{mol}$، أي \qty{29.6}{mg}. ولا
تتوقف النتيجة إلا على الشحنة، المقيسة من تيار وزمن، وعلى
ثابت فاراداي: فلا حاجة إلى أي محلول قياسي، بشرط أن يكون التحليل الكهربائي
تامًا ومردوده الفارادي 100\,\%.
\end{solution}

\begin{solution}{exo:b3:electrode-kinetics:11}
$\sqrt{DFv/RT} = \sqrt{\num{1.0e-5} \times 38.92 \times 0.1} = \qty{6.24e-3}{cm.s^{-1}}$: $\Lambda = 1.6$، \omterm{def:b3:electrode-kinetics:reversibility}{شبه عكوسة}.
وعند \qty{10}{V/s}، $\Lambda = 0.16$، لا تزال \omterm{def:b3:electrode-kinetics:reversibility}{شبه عكوسة} لكنها أبعد عن الحد
العكوس؛ ولن تكون الموجة عكوسة إلا دون نحو \qty{1}{mV/s}.
\end{solution}

\begin{solution}{exo:b3:electrode-kinetics:12}
$s/b = \qty{0.0838}{mM}$ و$b^2S_{xx} = \qty{203.5}{\micro A^2}$. عند $\bar y$: $0.0838\sqrt{1 + 1/7} = \qty{0.090}{mM}$.
وعند \qty{18.0}{\micro A}: $0.0838\sqrt{1.143 + 83.0/203.5} = \qty{0.104}{mM}$. ومن ثلاث قراءات عند $\bar y$:
$0.0838\sqrt{1/3 + 1/7} = \qty{0.058}{mM}$.
\end{solution}

\begin{solution}{pb:b3:electrode-kinetics:1}
\textbf{1.} $\ce{C6H12O6 -> C6H10O6 + 2 H+ + 2 e-}$ و$\ce{[Fe(CN)6]^3- + e- -> [Fe(CN)6]^4-}$؛
والتفاعل الإجمالي
\[
  \ce{C6H12O6 + 2 [Fe(CN)6]^3- -> C6H10O6 + 2 [Fe(CN)6]^4- + 2 H+} .
\]
وعند القطب، $\ce{[Fe(CN)6]^4- -> [Fe(CN)6]^3- + e-}$.
\textbf{2.} اثنان.
\textbf{3.} الأكسجين المذاب في الدم قليل ومتغير، وناتجه،
بيروكسيد الهيدروجين، لا يتأكسد إلا عند جهد مرتفع تتفاعل عنده أنواع
أخرى كثيرة؛ أما الوسيط فموجود بفائض معلوم.
\textbf{4.} كي يكون تركيزه عند السطح معدومًا: فلا يحدّ التيار عندئذ
إلا الانتشار، ولا يتأثر بالتغيرات الصغيرة في الجهد.
\textbf{5.} يتناقص التيار وفق $t^{-1/2}$؛ فيجب أخذ القراءات في الزمن المستعمل في
المعايرة.
\textbf{6.} كوتريل: $i\sqrt t = 42.0$، 41.4، 41.7، 41.6، 41.8: ثابت في حدود 1\,\%.
\textbf{7.} $D = \pi(\text{الميل}/FAc)^2 = \pi(\num{41.8e-6}/(96\,485 \times 0.0300 \times \num{1.00e-5}))^2 = \qty{6.55e-6}{cm^2.s^{-1}}$.
\textbf{8.} $\sqrt{\pi \times \num{6.55e-6} \times 5} = \qty{0.010}{cm} = \qty{100}{\micro m}$: فالاستنزاف يبلغ أعلى
الطبقة عند نحو \qty{5}{s}؛ والقراءات اللاحقة ستقع دون قانون كوتريل.
\textbf{9.} النصف: \qty{9.35}{\micro A}.
\textbf{10.} تيار الخلفية: أكسدة المتداخلات والشوائب، 
والشحن، وسداسي سيانوفيرات(II) المتبقي في الوسيط نفسه.
\textbf{11.} نعم: البواقي، التي لا تتجاوز \qty{0.095}{\micro A}، تتشتت دون اتجاه. وعند تراكيز الغلوكوز المرتفعة يتشبّع الوسيط، الموجود
بكمية محدودة، أو \omterm{def:b3:complex-kinetics:enzyme}{الإنزيم}، فينحني المستقيم.
\textbf{12.} $(7.10 - 0.356)/0.9189 = \qty{7.34}{mM}$.
\textbf{13.} $0.0838\sqrt{1 + 1/7 + (7.10 - 8.89)^2/203.5} = 0.0838 \times 1.076 = \qty{0.090}{mM}$.
\textbf{14.} الحد $1/m$ (قراءة واحدة)؛ وتقلّصه القراءات المكررة، أو عدة أشرطة.
\textbf{15.} $3 \times 0.077/0.919 = \qty{0.25}{mM}$.
\textbf{16.} $7.34 \times 180.16 = \qty{1322}{mg.L^{-1}} = \qty{132}{mg.dL^{-1}}$.
\textbf{17.} يضيف تيار أكسدته الخاص: فيقرأ الجهاز قيمة أعلى مما ينبغي.
\textbf{18.} يُظهر \omterm{def:b3:electrode-kinetics:cv}{فولتاموغرام} الشريط دون غلوكوز موجة مصعدية
للأسكوربات حيث يتأكسد سداسي سيانوفيرات(II).
\textbf{19.} تتأكسد مواد أقل عند جهد أدنى.
\textbf{20.} يُطرح تيار الشاهد من قراءة القطب الإنزيمي
قبل تطبيق مستقيم المعايرة.
\textbf{21.} يكبر $D$ ببضعة في المئة لكل كلفن، والتيار وفق $\sqrt D$: فالأجهزة تقيس
درجة الحرارة وتصحّح أثرها.
\textbf{22.} تغيّر اللزوجةُ ونسبةُ الكريات الحمراء في الدم $D$ والتيار؛
والمحاليل القياسية في وسط مماثل تلغي هذه الآثار.
\textbf{23.} تضاف تغيرات ما بين الأشرطة (مساحة القطب، وحمولة \omterm{def:b3:complex-kinetics:enzyme}{الإنزيم} والوسيط)،
ودرجة الحرارة وتركيب الدم إلى ارتياب المعايرة نفسها.
\textbf{24.} \textbf{$c(\text{الغلوكوز}) = 7.34 \pm \qty{0.09}{mmol.L^{-1}}$} (ارتياب معياري)، أي نحو \qty{132}{mg.dL^{-1}}.
\end{solution}
